\chapter {Policy-Based Deep Reinforcement Learning - From REINFORCE to A2C}
\label{chap:a2c}

This chapter develops the second path identified in Section~\ref{sec:policyoptimization}. Instead of approximating $q_*$ and deriving the policy via $\argmax$, the policy is directly parameterized as $\pi(a|s, \theta)$ and $\theta$ is optimized by gradient ascent on a scalar performance measure. A value function may still be used to evaluate and improve the policy, but it is not used for action selection.

\section{The Policy Gradient Theorem}
Besides optimizing the policy directly, parameterized policy methods offer two further advantages over action-value methods. First, a soft-max policy $\pi(a|s, \theta)$ \citep[p.~322]{suttonReinforcementLearningIntroduction2020} assigns a learned probability to every action, allowing the agent to discover that a mixture of actions outperforms any single greedy choice. An $\epsilon$-greedy policy cannot do this because it commits to one action with probability 1 - $\epsilon$ and distributes the remaining $\epsilon$ uniformly, regardless of what the agent has learned about non-greedy actions. Second, with a continuous policy parameterization, small changes in $\theta$ produce small smooth changes in the action probabilities. With $\epsilon$-greedy action selection over action values, an arbitrarily small change in the estimated values can cause the $\argmax$ to switch from one action to another, producing a discontinuous change in the policy \citep[p.~324]{suttonReinforcementLearningIntroduction2020}.

For the episodic case, such as Pong, the policy parameter $\theta$ is optimized by a gradient ascent on a scalar performance measure defined as the value of the start state:
\begin{equation}
	J(\theta) \doteq v_{\pi_{\theta}}(s_0),
\end{equation}
where $v_{\pi_{\theta}}$ is the true value function for $\pi_{\theta}$ determined by $\theta$, and $s_0$ is the start state. The gradient $\nabla J(\theta)$, however, is not straightforward to compute since changing $\theta$ changes both the action selection and the state distribution that the policy induces \citep[p.~324]{suttonReinforcementLearningIntroduction2020}. Computing the effect of the policy on the state distribution requires the environment dynamics, which are typically unknown \citep[p.~325]{suttonReinforcementLearningIntroduction2020}. A solution to this is provided by the policy gradient theorem:
\begin{equation}\label{eq:policy_gradient_theorem}
	\nabla J(\theta) \propto \sum_s \mu(s) \sum_a q_\pi(s, a)\,\nabla \pi(a\,|\,s, \theta),
\end{equation}
where $\mu(s)$ is the on-policy state distribution that appears only as a weighting and does not need to be differentiated. By sampling transitions through interaction with the environment, the agent visits states according to $\mu(s)$ naturally, which allows the gradient to be estimated without knowledge of the environment dynamics.

\section{REINFORCE and the Variance Problem}
The policy gradient theorem contains $ q_\pi(s, a)$ which cannot be computed directly since it is the expected return for taking action $a$ in state $s$ and following $\pi$ thereafter. As can be seen in \eqref{eq:action_value_function}, a single sample return $G_t$ is an unbiased estimate by definition.

\section{Actor-Critic Methods}

\section{The Advantage Function and A2C}